\documentclass[11pt]{article}
\usepackage[a4paper,margin=27mm]{geometry}
\usepackage[T1]{fontenc}
\usepackage{lmodern,amsmath,amssymb,amsthm,mathtools}
\usepackage{microtype,booktabs,longtable,array,xurl}
\usepackage[hidelinks]{hyperref}
\newtheorem{theorem}{Theorem}[section]
\newtheorem{lemma}[theorem]{Lemma}
\newtheorem{corollary}[theorem]{Corollary}
\newtheorem{proposition}[theorem]{Proposition}
\theoremstyle{definition}
\newtheorem{definition}[theorem]{Definition}
\theoremstyle{remark}
\newtheorem{remark}[theorem]{Remark}
\newcommand{\R}{\mathbb R}
\newcommand{\N}{\mathbb N}
\newcommand{\Z}{\mathbb Z}
\newcommand{\dens}{\operatorname{dens}}
\newcommand{\sgn}{\operatorname{sgn}}
\newcommand{\dist}{\operatorname{dist}}
\newcommand{\code}[1]{\texttt{\small\detokenize{#1}}}
\DeclareMathOperator{\Int}{int}
\title{Large closed sets avoiding a continuum of power asymptotics}
\author{}
\date{Research manuscript, 7 October 2026}
\begin{document}
\maketitle
\begin{abstract}
For every prescribed nonempty countable family of positive configurations with
bounded gaps in their occupied dyadic logarithmic bins, and every
$0<\varepsilon<1$, we construct a closed nowhere dense one-periodic set whose
measure in every unit interval is strictly greater than $1-\varepsilon$.
It simultaneously misses infinitely many distinct values in every tail of every
map with a nonzero leading power and a power-controlled higher order remainder.
The positive leading exponent, translation, coefficient and remainder map all
vary after the avoiding set has been chosen. In particular a compact subset of
$[0,1]$ of measure greater than $1-\varepsilon$ avoids every affine null geometric
progression, with every ratio in $(0,1)$, in every tail. The proof uses finite
random routing on nested periodic grids, exact representatives of a continuum
of parameters including boundary strata, and open repair of the closed set of
missed centers. A logarithmic window schedule makes the finer grid compatible
with positive power remainders. The traditional proof is given in full; a
separate correspondence records its relation to frozen Lean proofs.
\end{abstract}

\section{Statement, quantifiers and context}\label{sec:statement}
Let $\mu$ be Lebesgue measure. All real powers below have positive bases.
A set $S\subset\R$ is one-periodic if $S+1=S$; for a measurable such set write
$\dens S=\mu(S\cap[0,1))$. The integral of a one-periodic indicator over any
interval of length one is the same: after reducing its left endpoint modulo
one, split the interval there and translate its second part by $-1$.
Endpoints have measure zero, so closed unit intervals give the same value.

\begin{definition}\label{def:syndetic}
For $A\subset(0,\infty)$ put
\[
 S(A)=\{j\geq1:A\cap(2^{-j-1},2^{-j}]\ne\varnothing\}.
\]
We say that $A$ is logarithmically syndetic near zero if there are integers
$G,J\geq1$ such that every $\{j,\ldots,j+G-1\}$ with $j\geq J$ meets $S(A)$.
\end{definition}
The constants $G,J$ may depend on the configuration. The sets $A$ need not be
measurable or countable. They have arbitrarily small positive members.

\begin{theorem}[Prescribed family and continuum of power asymptotics]
\label{thm:main}
Fix a nonempty countable index set $\Lambda$ and configurations
$(A_\ell)_{\ell\in\Lambda}$ that are logarithmically syndetic near zero.
For every $0<\varepsilon<1$ there is a closed nowhere dense one-periodic
$E\subset\R$ such that
\begin{equation}\label{eq:density}
 \mu(E\cap[x,x+1])>1-\varepsilon\qquad(x\in\R).
\end{equation}
The same $E$ has the following property for every $\ell\in\Lambda$,
$s>0$, $\alpha>0$, $y\in\R$, $c\ne0$, and finite $M\geq0$.
Let $\delta>0$ and let $f$ be defined on
$D=A_\ell\cap(0,\delta)$. Suppose that, for some $\rho_0>0$,
\begin{equation}\label{eq:remainder}
 |f(a)-y-ca^s|\leq M a^{s+\alpha}
 \quad(a\in D,\ 0<a<\rho_0).
\end{equation}
Then for every $\rho>0$ the set
\begin{equation}\label{eq:distinct}
 \{f(a):a\in D,\ a<\rho,\ f(a)\notin E\}
\end{equation}
is infinite. No continuity, measurability or injectivity of $f$ is assumed.
\end{theorem}
Thus the family precedes the choice of $E$, while all the parameters and maps
in \eqref{eq:remainder} follow it. The theorem does not choose one set for all
possible configurations. Infinitude refers to distinct output values, rather
than merely infinitely many input indices. The eventual hypothesis permits
partial functions on a tail without any extension regularity requirement.

\begin{corollary}[Compact version]\label{cor:compact}
Under the hypotheses of Theorem~\ref{thm:main} there is a compact nowhere dense
$F\subset[0,1]$ with $\mu(F)>1-\varepsilon$ for which \eqref{eq:distinct}, with
$E$ replaced by $F$, is infinite for all the stated maps and parameters.
\end{corollary}
\begin{corollary}[All ratios of affine null geometric progressions]
\label{cor:geometric}
For every $0<\varepsilon<1$ there is a compact nowhere dense $F\subset[0,1]$
with $\mu(F)>1-\varepsilon$ such that
\[
 \forall a\ne0\ \forall b\in\R\ \forall q\in(0,1)\ \forall N\in\N,
 \quad\{a q^n+b:n\geq N\}\setminus F\text{ is infinite}.
\]
In particular no tail of any such progression is contained in $F$.
\end{corollary}
The proof of the main theorem occupies Sections~\ref{sec:sample}--\ref{sec:finish}.
The corollaries are proved in Section~\ref{sec:finish}.

\subsection{Related work and the role of the normalization}
Burgin, Goldberg, Keleti, MacMahon and Wang~\cite{BGKMW} ask in their
Question~1 whether every positive measure measurable subset of $\R$ contains
some $(a q^n+b)_{n\geq1}$, with $a\ne0$ and $0<q<1$. Their Theorem~1.1
constructs a compact set with zero as a density point avoiding all nonconstant
\emph{untranslated} $(a q^n)_{n\geq1}$. The translation in their question is
therefore material. Corollary~\ref{cor:geometric} supplies a negative answer to
that question, with the stronger measure and tail conclusions stated above.
This is a logical comparison of the exact statements, not a claim that a
literature search certifies exclusive priority.

Kolountzakis and Papageorgiou~\cite{KP}, Theorem~4.1, use finitely many
representatives of transformation parameters obtained from grid-crossing
hyperplanes. Their work provides precedent for probabilistic grid avoidance
and polynomial parameter arrangements; its main setting concerns sequences
escaping to infinity under growth hypotheses. The planar sign-arrangement
argument here is elementary and is proved below. The power normalization
$q=2^{-s}$ and the dyadic coefficient normalization $|c|=2^k t$ are also
elementary reductions, not separate contributions.

The construction integrates the nested-grid routing and exceptional-center
repair recorded in the pinned geometric-case research source~\cite{Routing}
and in the nonlinear avoidance entry~\cite{Profiles}. The prescribed countable
logarithmic profile perspective of the latter motivates working in output
logarithmic windows. The former chooses its avoiding set for a fixed ratio;
the latter prescribes a countable profile family and permits positive
logarithmic upper Banach density. We impose the stronger bounded-gap
configuration hypothesis while allowing a continuum of leading exponents;
neither parameter scope nor configuration scope should be conflated.
The specific extension proved here handles a continuum
of exponents \emph{inside} each finite routing lemma: active input indices
depend on the exponent, exact affine signs determine their grid addresses,
and a coupled logarithmic-window and power-buffer schedule keeps the error
cost small. We give the conditioning and repair arguments anew, so no theorem
from those research sources is a premise of the traditional proof.

The power gain $\alpha>0$ is a substantive boundary. Feng, Lai and
Xiong~\cite{FLX}, Theorem~1.1, show that positive decreasing sequences with
$\limsup a_{n+N}/a_n<1$ for some $N$ embed into every positive measure set by a
bi-Lipschitz map differentiable at zero with derivative one. For a dyadic
sequence this gives a remainder $o(a)$, without a promised $O(a^{1+\alpha})$
rate. Accordingly Theorem~\ref{thm:main} asserts avoidance of its specified
power-controlled maps, not arbitrary differentiable, $C^1$, or flat germs.
The source-specific comparisons and their limitations are documented in the
supplement; no exhaustive novelty or current-status certification is asserted.

\section{The robust compact-parameter blocker}\label{sec:sample}
\begin{lemma}[Uniform robust blocker]\label{lem:blocker}
Fix a logarithmically syndetic $A\subset(0,\infty)$, real numbers
$0<s_0<s_1$, $\alpha_0>0$, integers $k\in\Z$, $q,h\geq1$, and $0<p<1/12$.
There is an open one-periodic $H\subset\R$ with $\dens H\leq6p$ such that
for every $x\in\R$, $s\in[s_0,s_1]$, $t\in[1,2]$, and arbitrary errors
$e:A\cap(0,2^{-h})\to\R$ satisfying
\[
 |e(a)|\leq q a^{s+\alpha_0},
\]
some $a\in A\cap(0,2^{-h})$ satisfies
\begin{equation}\label{eq:hit}
 x+t2^k a^s+e(a)\in H.
\end{equation}
The set $H$ is chosen before $x,s,t,e$.
\end{lemma}
The next six sections prove this lemma. All probability spaces used are finite.
The errors in its statement are deterministic arbitrary choices; no independence
assumption about them is involved.

\begin{lemma}[Separated sampling and open-window counts]\label{lem:sampling}
There are $a_n\in A\cap(0,2^{-h})$ decreasing to zero, and $B_0>0$, such that
with $z_n=-\log_2 a_n$,
\begin{equation}\label{eq:loggaps}
 3/s_0<z_{n+1}-z_n<B_0.
\end{equation}
Put $D=s_1B_0$ and $\eta=1/(2D)$. For every $s\in[s_0,s_1]$ the sequence
$w_n(s)=s z_n-k$ increases to infinity with successive gaps greater than $3$
and less than $D$. If $u\geq s_1z_1-k$ and $L\geq2D$, then
\begin{equation}\label{eq:counts}
 \eta L\leq\#\{n:u<w_n(s)<u+L\}\leq1+L/3.
\end{equation}
The inequalities apply at parameters where a member equals a window endpoint.
\end{lemma}
\begin{proof}
Choose an occupied bin $j_1>\max(J,4,h)$, and let
$m=\lceil3/s_0\rceil+2$. Having chosen $j_n$, syndeticity permits an occupied
$j_{n+1}\in\{j_n+m,\ldots,j_n+m+G-1\}$. Choose $a_n$ in its bin.
Then $z_n\in[j_n,j_n+1)$, so
$z_{n+1}-z_n>m-1>3/s_0$ and $z_{n+1}-z_n<m+G$.
Take $B_0=m+G$. These bounds also prove convergence and the assertions about
$w_n(s)$. There is a last member at or below $u$ and a first member at or
above $u+L$. If $r$ members lie strictly between these endpoints, the length
$L$ is covered by $r+1$ successive gaps, giving $L\leq(r+1)D$.
Thus $r\geq L/D-1\geq\eta L$. Separation greater than $3$ gives the upper
bound. Equality with either endpoint is deliberately excluded from $r$.
\end{proof}

\section{Routing windows and actual periodic grids}\label{sec:routing}
Fix the data of Lemma~\ref{lem:blocker} and the sample above. Choose a branching
integer $b\geq2$ such that $p\eta(b-1)\geq12$, and put
\begin{equation}\label{eq:rates}
 \lambda=p\eta(b-1)/2,\qquad \kappa=\lambda-4\log2>0.
\end{equation}
Here and subsequently unmarked logarithms are natural. Choose $d\geq1$ so that
$(1-2^{1-b})^d<p$. Let $K$ be the number of edges in the complete ordered
$b$-ary tree of depth $d$. Choose an integer $g\geq1$ with
\begin{equation}\label{eq:stablecost}
 K2^{3-g}<p.
\end{equation}
All of $b,d,K,g,\eta,\lambda,\kappa$ are fixed before choosing the late
position of the routing template.

Fix integers $U,L$ satisfying the standing conditions
\[
 U\geq4,\qquad U>|k|,\qquad U\geq s_1z_1-k,\qquad
 L\geq\max\{2,\lceil2D\rceil\}.
\]
These conditions apply throughout the routing and probability arguments below;
Section~\ref{sec:schedule} will impose a sufficiently late schedule on the same
parameters. For this base length $L$, define
\begin{equation}\label{eq:recurrence}
\begin{aligned}
 L_1&=L,\qquad & T_1&=bL+(b-1)g,\\
 L_r&=g+T_{r-1},\qquad &
 T_r&=b(L_r+g+T_{r-1})+(b-1)g\quad(r\geq2).
\end{aligned}
\end{equation}
The height of a vertex is its remaining distance to the leaves. Each of its
outgoing edges receives an integer window of length $L_r$ when its height is
$r$. Traverse edges in preorder, including default edges, placing $g$ unused
integer indices between consecutive edge windows. Start at this integer $U$.
For edge $e$ write its endpoints $u_e,v_e=u_e+L_r-1$, and write $v_e^*$ for
the last endpoint of that edge and its entire child subtree. Set $T=T_d$.

\begin{lemma}[Template spans]\label{lem:template}
The last endpoint of the template is $U+T-1$, every edge length is at least
$L$, and for an edge of length $L_r$,
\begin{equation}\label{eq:localspan}
 v_e^*-u_e+1\leq2L_r.
\end{equation}
For fixed $b,d,g$, the total span $T$ is affine in $L$ with positive finite
integer coefficients.
\end{lemma}
\begin{proof}
A height-one subtree has the $b$ windows of length $L$ separated by $b-1$
gaps. At height $r\geq2$ one edge followed by its child subtree has span
$L_r+g+T_{r-1}=2L_r$. There are $b$ such pieces and $b-1$ intervening gaps,
giving \eqref{eq:recurrence}. This also proves the endpoint formula by induction.
At height one the local span is only $L_1$. Finally
$T_r=2bT_{r-1}+(3b-1)g$ for $r\geq2$, proving the affine assertion.
\end{proof}

For every integer endpoint $v$ use the periodic grid of $N_v=2^{v+3}$ cells,
with address
\[
 J_v(x)=\lfloor N_vx\rfloor\pmod{N_v}.
\]
Its cells are left closed and right open. The grids are nested. At each
nondefault edge (child number $i<b$) put a table of independent fair selector
bits indexed by its endpoint-grid cells. At each leaf put an independent
Bernoulli-$p$ terminal table indexed by the endpoint grid of its incoming edge.
All entries of all tables are mutually independent. Route a point from the
root by taking the first nondefault child with selector bit one at its address,
or child $b$ if all the nondefault bits are zero. Include the point in $B$ when
its terminal bit is one. For each table outcome $B$ is a measurable periodic
union of cells of the single finest grid
\begin{equation}\label{eq:finest}
 N=2^{U+T+2}.
\end{equation}
For a fixed point, conditioning on every selector leaves its addressed terminal
entry Bernoulli-$p$. Finite summation followed by integration over $[0,1)$ gives
\begin{equation}\label{eq:expectedB}
 \mathbb E\dens B=p.
\end{equation}

An index $n$ is \emph{active} on edge $e$ at exponent $s$ precisely when
\begin{equation}\label{eq:activation}
 u_e<w_n(s)<v_e+1=u_e+L_r.
\end{equation}
These are fixed \emph{output} windows. The input indices supplying their tests
may vary with $s$. This feature avoids requiring one input window to separate
outputs at $s_1$ while preserving their earlier addresses at $s_0$.

Call a center $x$ stable if, for every noninitial window starting at
$u_e=v+g+1$, no boundary of the grid $N_v$ lies in
$(x,x+2^{-v-g}]$. This is a deterministic measurable periodic condition.

\begin{lemma}[Stability and address separation]\label{lem:geometry}
The unstable centers have density at most $K2^{3-g}<p$. Assume $U\geq4$.
For a stable $x$, any active point
$x+t2^ka_n^s$, $t\in[1,2]$, has all earlier preorder-grid addresses equal to
those of $x$. On its own edge grid its address differs from the center's,
and two distinct active indices on that edge have different addresses.
These distinctions persist on every finer grid.
\end{lemma}
\begin{proof}
For a preceding endpoint $v$, the bad-center intervals have total length at
most $N_v2^{-v-g}=2^{3-g}$ per period. Sum over at most $K$ predecessors.
An active offset satisfies
\begin{equation}\label{eq:offsetbounds}
 2^{-u_e-L_r}<t2^{k-sz_n}<2^{1-u_e}.
\end{equation}
If there is a predecessor $v$, the upper bound is $2^{-v-g}$, so stability
prevents crossing its boundary. All earlier grids divide that grid; hence their
addresses agree as well. The own grid has width $2^{-u_e-L_r-2}$, equal to
one quarter of the displayed lower offset bound. Different active logarithms have gap
greater than $3$, so their positive offsets have ratio greater than $8$.
Their difference is greater than $7\,2^{-u_e-L_r}$, also exceeding the cell
width. All offsets lie in an interval of length less than $1/8$ since $U\geq4$,
so the separations cannot be hidden by periodic wraparound. Nested finer cells
preserve the distinctions. Floor addresses at exact grid boundaries obey the
same left-closed convention throughout.
\end{proof}

\section{Finite probability with center exposure}\label{sec:probability}
Fix a stable center $x$. Expose its addressed selector entry in \emph{every}
selector table, not only the tables on its route. These entries determine its
route and partition the finite probability space into center-exposure atoms.
The probability that this route never takes a default edge is
\begin{equation}\label{eq:nodefault}
 (1-2^{1-b})^d<p.
\end{equation}
Indeed at each reached vertex the fresh $b-1$ exposed fair entries are all zero
with probability $2^{1-b}$, independently of the preceding route choices.

Consider an exposure atom whose first default is at a height-$r$ vertex.
Fix $(s,t)\in[s_0,s_1]\times[1,2]$. For each outgoing nondefault edge $e_i$
and each of its active indices, take a local test: its own selector bit times
the terminal bit obtained by routing the point within child subtree $i$.
By Lemma~\ref{lem:geometry} the point agrees with the center at strict ancestors
and at preceding siblings of this default vertex. They have the requisite
exposed zero bits at those siblings. If its own selector is one and its local
terminal bit is one, it belongs to the actual global set $B$.

\begin{lemma}[Fixed-parameter joint miss]\label{lem:joint}
Under the standing conditions of Section~\ref{sec:routing}, on each
first-default exposure atom, at fixed
$(s,t)$, the probability that all these local tests fail is at most
$\exp(-\lambda L_r)$. Equivalently, without dividing by the mass of the atom,
its intersection with the miss event has probability at most its own mass
times $\exp(-\lambda L_r)$.
\end{lemma}
\begin{proof}
The number $m_*$ of tests is at least $\eta(b-1)L_r$ by
Lemma~\ref{lem:sampling}, since $u_e\geq U\geq s_1z_1-k$ and
$L_r\geq L\geq\lceil2D\rceil$ on every outgoing edge.
Their own selector coordinates are distinct and avoid
every exposed coordinate: for one edge this follows from its own address
separation, and different edges have different tables. Thus the tested bits
$A_1,\ldots,A_{m_*}$ are free independent fair bits after center exposure.

Condition next on \emph{all} selector entries. Local terminal addresses are
pairwise distinct for every such assignment. Points in different child
subtrees or different leaves use different terminal tables. Points ending at
the same leaf started on the same child edge, where their own-grid addresses
were distinct, and keep distinct addresses at that leaf's finer grid. Hence
terminal independence gives conditional failure probability
$(1-p)^{\sum_j A_j}$. Averaging over the free own-selector bits yields exactly
$(1-p/2)^{m_*}$, irrespective of shared auxiliary selector reads or changes in
the terminal addresses. It is at most
$\exp(-pm_*/2)\leq\exp(-\lambda L_r)$.
Multiplying by exposure-atom masses and using finite sums proves the joint
form, including zero-mass atoms. No independence of entire routing paths is
assumed.
\end{proof}

\section{Exact representatives of the continuum}\label{sec:representatives}
We give a boundary-complete finite reduction at a fixed center and fixed
first-default vertex. It depends on the prescribed sample and actual grids,
and is chosen before the unexposed tables.

\begin{lemma}[Planar sign representatives]\label{lem:arrangement}
For $m$ affine real functions on $\R^2$, at most $2m^2+1$ different vectors
of signs in $\{-1,0,1\}^m$ are realized. Consequently any subset of the plane
admits a choice of at most $20(m+5)^2$ representatives realizing all its sign
vectors. All zero signs are retained.
\end{lemma}
\begin{proof}
Constant functions have fixed signs. Merge coincident zero lines of nonconstant
functions, keeping their individual oriented labels. With $r\leq m$ distinct
lines there are at most $r(r-1)/2$ intersection vertices. Each line is cut by
at most $r-1$ distinct intersection points into at most $r$ open edges, giving
at most $r^2$ edges. Insert the lines one by one; the $j$th line is cut into
at most $j$ pieces, each of which increases the number of open faces by at
most one. There are therefore at most $1+r(r+1)/2$ faces. These strata cover
the plane and have constant full sign vectors. Their total is at most
$2r^2+1$. Intersecting with an arbitrary prescribed subset introduces no new
sign vectors: choose one point for each vector realized in that subset.
Parallel lines, repeated intersections and constant zero functions only reduce
these counts. The larger displayed representative bound is convenient below.
\end{proof}

Let $L_r=\ell$ at the first-default vertex, and suppose $U>|k|$.
An index potentially active on edge $e$ must have
\[
 (u_e+k)/s_1\leq z_n\leq(u_e+\ell+k)/s_0.
\]
Since all $z_n>0$ have gaps greater than $3/s_0$, the number of candidates on
an edge is at most $\lceil(U+T+|k|)/3\rceil+1$. Thus the number of candidate
edge-index pairs is bounded by
\begin{equation}\label{eq:P}
 P=(b-1)\bigl(\lceil(U+T+|k|)/3\rceil+1\bigr).
\end{equation}
For each pair record two activation cuts $sz_n-k=u_e,u_e+\ell$.
Record also every lifted boundary of its \emph{local finest} grid, namely the
grid at $v_e^*$, whose offset from $x$ is
$\beta=j/N_{v_e^*}-x\in(0,2^{1-u_e})$. By
\eqref{eq:localspan} there are at most
\[
 1+N_{v_e^*}2^{1-u_e}\leq1+2^{2\ell+3}
\]
such boundaries. Boundary crossing is the equality
$t2^{k-sz_n}=\beta$. In coordinates $(s,v)=(s,\log t)$ it is the affine cut
\begin{equation}\label{eq:cut}
 v-sz_n\log2+k\log2-\log\beta=0.
\end{equation}
The sign equals the sign of the output's difference from that lifted grid
boundary, because logarithm is strictly increasing. The activation signs
and the signs in \eqref{eq:cut} therefore determine the vector that gives
\emph{none} for inactive pairs and the actual local finest-grid address for
active ones. That grid determines every table read in the pair's subtree.

There are at most $m=P(3+2^{2\ell+3})$ cuts. Apply
Lemma~\ref{lem:arrangement} to the closed rectangle
$[s_0,s_1]\times[0,\log2]$ and pull representatives back to $(s,t)$.
The resulting set has size at most
\begin{equation}\label{eq:repcount}
 20\bigl[P(3+2^{2\ell+3})+5\bigr]^2.
\end{equation}
At an activation equality that pair is inactive. At a grid equality it has
the floor key of the cell to the right. The sign zero records both conventions
exactly. Parameter endpoints, intersections, coincident cuts and lower
dimensional strata are covered by the same argument. In particular this is
not a numerical sampling of exponents, and the representatives work for
\emph{every} assignment of all remaining table entries.

\begin{lemma}[Continuum miss bound]\label{lem:entropy}
Under the same standing conditions, let $C$ be a first-default center-exposure
atom, and let $\mathcal M_C$ be the event that, at its first-default vertex,
all active local tests fail for some $(s,t)\in[s_0,s_1]\times[1,2]$. Then the
joint probability satisfies
\begin{equation}\label{eq:entropy}
 \mathbb P(C\cap\mathcal M_C)
 \leq\mathbb P(C)\,5120 b^2(U+T+|k|+2)^2 e^{-\kappa\ell}.
\end{equation}
\end{lemma}
\begin{proof}
Every miss has the same active-address vector as a retained representative,
and therefore the same local tests for every outcome. Use the finite union
bound, \eqref{eq:repcount}, and Lemma~\ref{lem:joint} at each representative.
The factor multiplying $\mathbb P(C)$ is
$20[P(3+2^{2\ell+3})+5]^2 e^{-\lambda\ell}$.
For $P\geq1$ and $\ell\geq2$,
\[
 P\leq b(U+T+|k|+2),\qquad
 P(3+2^{2\ell+3})+5\leq16P2^{2\ell}.
\]
Squaring and using \eqref{eq:rates} gives \eqref{eq:entropy}.
The factor depending on the absolute position $U$ has thus been retained.
Summing the joint estimates over atoms avoids division by any atom mass.
\end{proof}

\section{A coupled window and error-buffer schedule}\label{sec:schedule}
For integer $U$ define
\begin{equation}\label{eq:Lschedule}
 L(U)=\max\left\{g,\lceil2D\rceil,2,
 \left\lceil\frac{4\log(U+2)+\log(46080b^2/p)+1}{\kappa}\right\rceil\right\}.
\end{equation}
Use this length in \eqref{eq:recurrence}, writing the resulting span as $T(U)$.
By Lemma~\ref{lem:template}, $T(U)=O(\log U)$, with a finite constant depending
on the already fixed tree. In particular $T(U)/U\to0$.

\begin{lemma}[Simultaneous late schedule]\label{lem:schedule}
There are arbitrarily large integers $U$ for which
\begin{equation}\label{eq:late}
 U\geq4,\qquad U>|k|,\qquad U\geq s_1z_1-k,\qquad T(U)\leq U,
\end{equation}
every first-default continuum miss bound is less than $p$ times its atom mass,
and a common positive error radius $r_U$ has $4Nr_U<p$, where $N$ is the
actual finest grid \eqref{eq:finest}.
\end{lemma}
\begin{proof}
The conditions \eqref{eq:late} hold for all sufficiently large $U$.
For $\ell\geq L(U)$, $T(U)\leq U$, and $|k|<U$, the factor multiplying the
atom mass in \eqref{eq:entropy} is at most
\begin{align}
 46080b^2(U+1)^2e^{-\kappa L(U)}
 &\leq p e^{-1}\frac{(U+1)^2}{(U+2)^4}<p.\label{eq:smallentropy}
\end{align}
For any active input on any edge, $sz_n-k>U$ and $s\leq s_1$ give
$2^k a_n^s<2^{-U}$ and $z_n>(U+k)/s_1$. Since $U+k>0$, its allowed error
has the bound
\[
 q a_n^{s+\alpha_0}
 \leq q2^{-k}2^{-U}2^{-\alpha_0(U+k)/s_1}.
\]
Define
\begin{equation}\label{eq:radius}
 r_U=(q2^{-k}+1)2^{-U}2^{-\alpha_0(U+k)/s_1}>0.
\end{equation}
The allowed active error is strictly smaller than $r_U$. The doubled-buffer
cost on the \emph{actual} grid is
\begin{equation}\label{eq:buffercost}
 4Nr_U=2^{T(U)+4}(q2^{-k}+1)2^{-\alpha_0(U+k)/s_1}\longrightarrow0.
\end{equation}
Indeed its logarithm is $O(\log U)-\alpha_0(\log2)U/s_1+O(1)$, which tends
to $-\infty$. Increasing $U$ further enforces $4Nr_U<p$ while keeping all
preceding constraints. The order is noncircular: tree constants are fixed,
then $L(U),T(U)$ are explicit logarithmic functions, and all inequalities
hold eventually for those same functions. No practical bound on the size
of the finite probability space is asserted.
\end{proof}

\section{Closed missed centers and open repair}\label{sec:repair}
Fix a schedule from Lemma~\ref{lem:schedule}. For each outcome put
\[
 B_1=B+(-r_U,r_U),\qquad B_2=B+(-2r_U,2r_U).
\]
These are open periodic sets. A union of cells of an $N$-cell periodic grid
enlarged by radius $r$ gains density at most $2Nr$: enlarge each of its at
most $N$ cell intervals by $r$ at both ends and use subadditivity, allowing
periodic overlap. Thus \eqref{eq:expectedB} and \eqref{eq:buffercost} give
\begin{equation}\label{eq:expectedbuffer}
 \mathbb E\dens B_2\leq p+4Nr_U<2p.
\end{equation}

Only finitely many indices can be active on some template edge for some
$s\in[s_0,s_1]$, since their logarithms lie in
$[(U+k)/s_1,(U+T+k)/s_0]$ and are separated. Let $I$ contain precisely these
potential indices. For $n\in I$ let $S_n$ be the relatively open subset of
$[s_0,s_1]$ where $w_n(s)$ belongs to the union of the open template windows.
Define $R$ to be the set of centers for which some $(s,t)$ in the compact
parameter rectangle has
\begin{equation}\label{eq:missed}
 s\notin S_n\quad\hbox{or}\quad x+t2^k a_n^s\notin B_1
 \qquad\text{for every }n\in I.
\end{equation}

\begin{lemma}[Closedness and expected density of missed centers]
\label{lem:closed}
For each outcome $R$ is closed and one-periodic, and
$\mathbb E\dens R\leq3p$.
\end{lemma}
\begin{proof}
For each $n$, the first alternative in \eqref{eq:missed} is closed because
$S_n$ is relatively open. The second is closed by continuity of
$(x,s,t)\mapsto x+t2^ka_n^s$ and openness of $B_1$. Their union and the finite
intersection over $I$ are closed. Projection along a compact parameter space
preserves closedness: from a convergent sequence of projected centers choose
parameter witnesses, pass to a convergent subsequence in the rectangle, and
use closedness. Translation of $x$ by an integer preserves all conditions.
Thus $R$ is measurable as well.

At a stable center, membership in $R$ implies a continuum local miss at its
first-default vertex, unless the center route has no default. This follows
because a successful local test gives membership in $B\subset B_1$.
Sum \eqref{eq:smallentropy} over the first-default atoms and add
\eqref{eq:nodefault}; the pointwise probability is at most $2p$.
At an unstable center it is at most one. Finite summation and integration
over a period, together with \eqref{eq:stablecost}, give
\[
 \mathbb E\dens R
 \leq2p\mu(\text{stable}\cap[0,1))+
       \mu(\text{unstable}\cap[0,1))\leq3p.
\]
All exchanges are finite sums of measurable indicators. No measurable choice
of continuum representatives as a function of $x$ is needed for this
pointwise probability estimate.
\end{proof}

\begin{lemma}[Deterministic repair]\label{lem:repair}
There is a table outcome and an open periodic $V\supset R$ for which
$H=B_2\cup V$ has $\dens H\leq6p$ and satisfies \eqref{eq:hit} for every
center and every permitted error family.
\end{lemma}
\begin{proof}
By \eqref{eq:expectedbuffer} and Lemma~\ref{lem:closed}, some outcome has
$\dens B_2+\dens R\leq5p$. To justify the existence directly, a finite
weighted average with nonnegative weights summing to one cannot be smaller
than every positive-weight value; zero-weight outcomes cause no problem.

For this closed periodic $R$, the open thickenings
$V_j=\{x:\dist(x,R)<1/j\}$ decrease to $R$, are periodic, and their measures
restricted to $[0,1)$ decrease to $\dens R$ by continuity from above in a
finite measure space. (If $R$ is empty use $V=\varnothing$.) Choose $j$ with
$\dens V_j<\dens R+p$ and set $V=V_j$. Consequently
$\dens(B_2\cup V)<6p$.

Fix $x,s,t,e$. If $x\notin R$, the negation of \eqref{eq:missed} at these
parameters supplies a finite active input with unperturbed point in $B_1$.
Its error has absolute value at most $r_U$, so its perturbed point belongs
to $B_2$: a distance strictly less than $r_U$ from $B$ remains strictly less
than $2r_U$ after an error of size at most $r_U$. Equality of an allowed
error does not lose an open-buffer endpoint.

If $x\in R\subset V$, use the full sampled sequence. Its perturbed offsets
tend to zero; for $a_n<1$,
\[
 |t2^k a_n^s+e(a_n)|
 \leq2^{k+1}a_n^{s_0}+q a_n^{s_0+\alpha_0}\longrightarrow0.
\]
Since $V$ is open and contains $x$, a sufficiently late perturbed point lies
in $V$. This index need not lie in $I$ or be active on a finite window.
Every sampled input belongs to $A\cap(0,2^{-h})$. Thus every center is
repaired, including all exceptions, and Lemma~\ref{lem:blocker} follows.
\end{proof}

\section{Countable exhaustion and distinct tail values}\label{sec:finish}
\begin{proof}[Proof of Theorem~\ref{thm:main}]
The set of tuples
\[
 \tau=(\ell,N,j,k,q,h),\qquad
 \ell\in\Lambda,\quad N\geq2,\quad j,q,h\geq1,\quad k\in\Z,
\]
is countable. Choose positive $p_\tau<1/12$ such that
$12\sum_\tau p_\tau<\varepsilon$; for example, using an injection of the
tuples into $\N$, assign $p_\tau=\varepsilon2^{-n(\tau)-1}/24$.
For each tuple apply Lemma~\ref{lem:blocker} to $A_\ell$ with
$[s_0,s_1]=[1/N,N]$ and $\alpha_0=1/j$. Let $H_\tau$ be its blocker and put
\[
 \mathcal U=\bigcup_\tau(H_\tau\cup(-H_\tau)),\qquad E=\R\setminus\mathcal U.
\]
This is a countable union of open periodic sets with
$\dens\mathcal U\leq12\sum_\tau p_\tau<\varepsilon$.
Thus $E$ is closed and periodic, and the unit-interval identity in
Section~\ref{sec:statement} gives the strict inequality \eqref{eq:density}.

Now fix the configuration, map, and parameters in \eqref{eq:remainder}.
Choose $N$ with $s\in[1/N,N]$, $j$ with $1/j\leq\alpha$, and an integer
$q\geq\max(1,M)$. For $a<1$ the remainder bound is at most
$q a^{s+1/j}$. If $c>0$ choose $k\in\Z$ and $t\in[1,2)$ with $c=2^kt$.
For every sufficiently large $h$, $2^{-h}<\min(1,\delta,\rho_0)$.
Apply the corresponding blocker to $x=y$ and the errors
$e(a)=f(a)-y-ca^s$ on this tail. It gives $a_h\in A_\ell$, $a_h<2^{-h}$,
with $f(a_h)\in H_\tau\subset\mathcal U$.
If $c<0$, apply the same argument to $-f$, centered at $-y$, with coefficient
$-c>0$ and errors $-e(a)$; then $f(a_h)\in-H_\tau\subset\mathcal U$.
This uses only countably many budgets even though $s,\alpha,c,y,f$ range
after the set is chosen.

Since $M a^\alpha\to0$ and $c\ne0$, sufficiently small inputs satisfy
\begin{equation}\label{eq:twosided}
 \tfrac12|c|a^s\leq |f(a)-y|\leq\tfrac32|c|a^s.
\end{equation}
For any given $\rho>0$, choose hits above with $h$ arbitrarily large and
$a_h<\rho$. Their values lie outside $E$, differ from $y$, and tend to $y$.
They cannot form a finite set: a finite set avoiding $y$ has a positive
minimum distance from $y$, contrary to the upper bound in
\eqref{eq:twosided}. This proves infinitude of distinct values in every tail,
without assuming distinctness of chosen indices or injectivity of $f$.

Finally $E$ has empty interior. If it contained an open interval, choose a
point $y$ strictly inside that interval and any configuration in the nonempty
family. The map $f(a)=y+a$ on a sufficiently small tail would take all its
values inside the interval. It has $s=1,c=1,M=0$ and any $\alpha>0$, in
contradiction with \eqref{eq:distinct}. A closed set with empty interior is
nowhere dense. This completes the theorem.
\end{proof}

\begin{proof}[Proof of Corollary~\ref{cor:compact}]
Take $F=E\cap[0,1]$. It is compact, has measure strictly greater than
$1-\varepsilon$, and has empty interior because it is a subset of $E$.
Every value outside $E$ is outside $F$, so each infinite tail conclusion
passes to $F$.
\end{proof}

\begin{proof}[Proof of Corollary~\ref{cor:geometric}]
Use the singleton family $A=\{2^{-n}:n\geq1\}$. Every bin in
Definition~\ref{def:syndetic} contains its upper endpoint, so $G=J=1$ works.
For $q\in(0,1)$ set $s=-\log_2q>0$, $c=a$, and $y=b$.
The map $f(u)=b+au^s$ has zero remainder and gives
$f(2^{-n})=b+a q^n$. A sufficiently small input tail forces $n\geq N$.
Corollary~\ref{cor:compact} gives infinitely many distinct misses, simultaneously
for all the indicated parameters. In the separately frozen geometric proof,
the endpoint asks only for at least one miss in every index tail; the present
statement implies that endpoint and contains it as a corollary, without
counting it as another independent contribution.
\end{proof}

\section{Boundaries and an obstruction to this mechanism}\label{sec:scope}
The theorem requires $s>0$, $\alpha>0$, and $c\ne0$. A zero leading coefficient
does not force nonconstant values, and positive $\alpha$ is what makes
\eqref{eq:buffercost} tend to zero with a logarithmic total span. The proof
does not extend by merely replacing this gain with an arbitrary $o(a^s)$
bound. The compact parameter rectangle includes its entire boundary; activations
are open, while grid cells are left closed. These choices play different roles.

Here are two exact finite controls illustrating those roles. They supplement,
but do not replace, the infinite proof.
For inputs $a_1=1/2,a_2=1/4$, $s,t\in[1,2]$, and holes
$H=(1/10,1)+\Z$, arbitrary errors of size at most $r<7/160$ cannot keep both
outputs outside $H$, whatever the center $x$. Indeed with $u=2^{-s}$ their
positive unperturbed difference is $t(u-u^2)\in[3/16,1/2]$, and after errors
the difference lies strictly between $1/10$ and $9/10$. Two points in the
closed complementary gaps cannot have such a difference. At $r=7/160$ take
$s=2,t=1,x=143/160$ and errors $(-7/160,7/160)$: the perturbed points are
$11/10$ and $1$, both missed endpoints. For $s<2$ or $t>1$ the same worst
error bound still succeeds. Deleting the parameter boundary therefore loses
this sharp failure.

Closed activation would also invalidate the closed-projection argument.
For one input $a=1/2$, $x=0,t=1$, ambient exponent interval $[1/2,5/2]$,
closed activation $s\in[1,2]$ and holes
$(2/5,3/5)$, the condition ``inactive or missed'' holds along $s<1$ but fails
at the limit $s=1$. With open activation $(1,2)$ the limit is inactive and
the relevant union is closed, as used in Lemma~\ref{lem:closed}.

\begin{proposition}[Positive upper Banach density does not suffice for the
window-counting step]\label{prop:obstruction}
There is a set $S\subset\N$ of upper Banach density one such that, for
$A=\{2^{-j}:j\in S\}$ and each fixed finite $C>0$, every sufficiently large
$U$ and every $0<T\leq C\log U$ admit an $s\in[1,2]$ for which
\[
 \{j\in S:U<sj<U+T\}=\varnothing.
\]
This is an obstruction to the present uniform output-window mechanism, not
a disproof of an avoidance theorem for that larger class of configurations.
\end{proposition}
\begin{proof}
Put $R_m=2^{2^m}$ and
$S=\bigcup_{m\geq3}\{R_m,R_m+1,\ldots,R_m+m\}$. Its blocks have arbitrary
length, so for every integer interval length the maximum occupied proportion
is one, giving upper Banach density one. An active $j$ for $s\in[1,2]$ must
lie in $[U/2,U+T]$. For large $U$ we have $T\leq U/4$. Any block meeting that
interval has $R_m\leq5U/4$, hence $m=O(\log\log U)$, and
$R_m+m\geq U/2$, hence eventually $R_m\geq U/3$.
There is at most one such block: consecutive starts obey $R_{m+1}=R_m^2$,
so two starts in $[U/3,5U/4]$ would require $U^2/9\leq5U/4$.
If no block meets the interval, any $s$ works. Otherwise the exponents that
activate any member of that block lie in
\[
 \left(\frac{U}{R_m+m},\frac{U+T}{R_m}\right).
\]
The length is at most $3T/U+9m/U$, tending to zero uniformly for the stated
$T$. It is eventually less than one and cannot cover $[1,2]$. Choose an
exponent outside it.
\end{proof}

\appendix
\section{Formal correspondence and evidence status}\label{app:formal}
The formal sources and their dependency locks are preserved as two immutable
archives in the supplement. The stronger endpoint is
\code{ContinuumRemainder.continuum_power_target}; its compact endpoint is
\code{ContinuumRemainder.compact_power_avoidance}. The earlier endpoint is
\code{ContinuumGeometric.geometric_main_target}. The formal definitions use
total functions with an eventual bound. A function on
$A\cap(0,\delta)$ can be extended arbitrarily outside that domain; the exact
tail value sets below $\delta$ are unchanged. Conversely restriction of a
total function gives the partial-domain statement. The corresponding Lean
declarations are listed in the table. This is equivalence of the tail
quantifiers, not a continuity assertion about an extension.

The 18 stronger proof modules reuse 45 geometric modules byte for byte.
Two additional aggregate modules account for the independent audit's reported
65-module count. The complete machine-readable declaration inventory includes
source locations and hashes; the module import graph is provided separately.
Those import edges record dependencies between modules, not an asserted exact
constant-level kernel dependency graph. Every key proved step of the main
traditional argument is matched below to concrete formal declarations.
The elementary normalization in Corollary~\ref{cor:geometric} and the scope
observations are derived here; the earlier geometric endpoint is already
separately formalized. Proposition~\ref{prop:obstruction} is a traditional
scope result, not claimed to be a Lean theorem in these frozen modules.

The delivered normal and optimized harness reports both record successful
fresh compilation, expanded-target and partial-tail probes, axiom audits and
six rejected false controls. Their hashes agree. The previous archive's file
\code{evidence/KernelReplay.lean}, however, contains only an import; it does
\emph{not} implement explicit declaration replay into an empty kernel.
Consequently its original README's stronger description of that file is not
used as evidence here. The independent audit subsequently reported two
recompilations of 65 modules and true empty-environment, trust-zero kernel
replay of the closure of 1,445 safe declarations (35,620 declarations total).
These independent counts are attributed to that supplied audit report; its
underlying replay files were not included in the materials available for this
manuscript. The supplement distinguishes that report from locally available
reproducible harness evidence. Neither frozen proof archive has been modified.
The mathematical sources need no change on account of this evidence-description
correction. No external human referee or professional mathematician has been
represented as having reviewed this manuscript.

The fixed environment is Lean 4.34.1 and mathlib commit
\nolinkurl{d13f23b723b8a846827a245b89c10fc7d3f11612}; all transitive revisions are
retained in the supplied lockfile. The recorded axiom closures use only
\code{propext}, \code{Classical.choice}, and \code{Quot.sound}. The paper's
traditional proof is independently readable and does not appeal to a
proof-assistant theorem as an unproved mathematical premise.

\small
\begin{longtable}{@{}p{0.27\textwidth}p{0.68\textwidth}@{}}
\caption{Traditional proof steps and exact frozen Lean declarations.}\label{tab:correspondence}\\
\toprule Paper step & Lean declarations and source modules \\
\midrule\endfirsthead
\toprule Paper step & Lean declarations and source modules \\
\midrule\endhead
\bottomrule\endfoot
\ref{def:syndetic}. Logarithmic bins, per-configuration bounded gaps and exact endpoint specification & \nolinkurl{ContinuumRemainder.OccupiedBin}\par \nolinkurl{ContinuumRemainder.LogSyndetic}\par \nolinkurl{ContinuumRemainder.ContinuumPowerTarget}\par \nolinkurl{ContinuumRemainder.RobustCompactBlockerSpec}\par\smallskip Modules: \nolinkurl{ContinuumRemainder.Specification} \\\addlinespace
\ref{lem:sampling}. Separated sampling and strict output-window counts & \nolinkurl{ContinuumRemainder.LogSyndetic.exists_sampledLogConfiguration}\par \nolinkurl{ContinuumRemainder.open_window_count_bounds}\par \nolinkurl{ContinuumRemainder.SampledLogConfiguration.activeLabels_count_bounds}\par\smallskip Modules: \nolinkurl{ContinuumRemainder.Sampling}; \nolinkurl{ContinuumRemainder.Counting} \\\addlinespace
\ref{lem:template}. Local span, finest endpoints and affine total span & \nolinkurl{ContinuumGeometric.RoutingTemplate.blockSpan_le_twice_length}\par \nolinkurl{ContinuumGeometric.RoutingTemplate.span_affine}\par \nolinkurl{ContinuumGeometric.RoutingTemplate.edgeEnd_within_total}\par\smallskip Modules: \nolinkurl{ContinuumGeometric.RoutingTemplate}; \nolinkurl{ContinuumGeometric.RoutingTreeBounds} \\\addlinespace
\ref{lem:geometry}. Stable earlier addresses and distinct own-grid keys & \nolinkurl{ContinuumRemainder.log_stable_earlier_address_agreement}\par \nolinkurl{ContinuumRemainder.log_gridAddress_ne_center}\par \nolinkurl{ContinuumRemainder.log_gridAddress_ne_of_output_gap}\par \nolinkurl{ContinuumGeometric.unitDensity_routingUnstableCenters_le}\par\smallskip Modules: \nolinkurl{ContinuumRemainder.LogGeometry}; \nolinkurl{ContinuumGeometric.RoutingStableMeasure} \\\addlinespace
\ref{sec:routing}. Actual finite global grid and Bernoulli terminal expectation & \nolinkurl{ContinuumGeometric.actual_routedSet_finite_grid}\par \nolinkurl{ContinuumGeometric.actual_routedSet_expected_density}\par\smallskip Modules: \nolinkurl{ContinuumGeometric.RoutingGlobalGrid}; \nolinkurl{ContinuumGeometric.FiniteRoutingProbability} \\\addlinespace
\ref{sec:probability}. Full center exposure and exact no-default probability & \nolinkurl{ContinuumGeometric.actual_center_atom_probability}\par \nolinkurl{ContinuumGeometric.actual_center_atoms_sum_one}\par \nolinkurl{ContinuumGeometric.actual_routeHasNoDefault_probability}\par\smallskip Modules: \nolinkurl{ContinuumGeometric.FiniteRoutingProbability}; \nolinkurl{ContinuumGeometric.NoDefaultProbability} \\\addlinespace
\ref{lem:joint}. Joint atom miss identity without division or path independence & \nolinkurl{ContinuumGeometric.joint_center_atom_all_miss}\par \nolinkurl{ContinuumRemainder.log_fixed_parameter_joint_miss}\par\smallskip Modules: \nolinkurl{ContinuumGeometric.FiniteRoutingProbability}; \nolinkurl{ContinuumRemainder.LogRoutingProbability} \\\addlinespace
\ref{lem:arrangement}. Boundary-complete planar affine sign arrangement & \nolinkurl{ContinuumGeometric.realizedLinePatterns_card_le}\par \nolinkurl{ContinuumGeometric.realizedPlanePatterns_card_le}\par \nolinkurl{ContinuumGeometric.arrangementRepresentativeBound}\par\smallskip Modules: \nolinkurl{ContinuumGeometric.LineSignBound}; \nolinkurl{ContinuumGeometric.Planar} \\\addlinespace
\ref{sec:representatives}. Candidate count, activation signs and exact grid crossing signs & \nolinkurl{ContinuumRemainder.SampledLogConfiguration.potentialPairs_card_le}\par \nolinkurl{ContinuumRemainder.log_activation_cut_signs}\par \nolinkurl{ContinuumRemainder.log_grid_boundary_sign}\par \nolinkurl{ContinuumRemainder.log_grid_key_eq_of_log_signs}\par\smallskip Modules: \nolinkurl{ContinuumRemainder.Counting}; \nolinkurl{ContinuumRemainder.LogSignatures} \\\addlinespace
\ref{lem:entropy}. Representatives uniform in all outcomes and the continuum union bound & \nolinkurl{ContinuumRemainder.log_local_grid_representatives_uniform}\par \nolinkurl{ContinuumRemainder.log_representatives_entropy_bound}\par \nolinkurl{ContinuumRemainder.log_continuum_joint_miss_bound}\par \nolinkurl{ContinuumRemainder.sample_vertex_continuum_joint_miss_bound}\par\smallskip Modules: \nolinkurl{ContinuumRemainder.LogSignatures}; \nolinkurl{ContinuumRemainder.LogRoutingProbability}; \nolinkurl{ContinuumRemainder.SampleLocalProbability} \\\addlinespace
\ref{lem:schedule}. Logarithmic span, vanishing actual buffer cost and simultaneous late schedule & \nolinkurl{ContinuumRemainder.scheduledSpan_isBigO_log}\par \nolinkurl{ContinuumRemainder.sublinear_buffer_cost_tendsto_zero}\par \nolinkurl{ContinuumRemainder.exists_late_robust_schedule}\par \nolinkurl{ContinuumRemainder.exists_sample_routing_schedule_late}\par\smallskip Modules: \nolinkurl{ContinuumRemainder.ErrorSchedule}; \nolinkurl{ContinuumRemainder.SampleSchedule} \\\addlinespace
\ref{sec:repair}. Active error domination and strict open double buffers & \nolinkurl{ContinuumRemainder.active_power_error_lt_radius}\par \nolinkurl{ContinuumRemainder.active_error_in_outer_buffer}\par \nolinkurl{ContinuumGeometric.unitDensity_periodicGridSet_double_buffer_le}\par\smallskip Modules: \nolinkurl{ContinuumRemainder.ErrorDomination}; \nolinkurl{ContinuumGeometric.PeriodicRepair} \\\addlinespace
\ref{lem:closed}. Closed missed centers, compact projection and stable probability & \nolinkurl{ContinuumRemainder.isClosed_sampledMissedCenters}\par \nolinkurl{ContinuumRemainder.sampled_stable_missed_center_probability_le}\par \nolinkurl{ContinuumRemainder.sampled_stable_bad_density_le}\par\smallskip Modules: \nolinkurl{ContinuumRemainder.RobustRepair}; \nolinkurl{ContinuumRemainder.SampleStableProbability} \\\addlinespace
\ref{lem:repair}. Open cover, every-center convergence and outcome assembly & \nolinkurl{ContinuumGeometric.closed_onePeriodic_open_cover}\par \nolinkurl{ContinuumRemainder.perturbed_power_tendsto}\par \nolinkurl{ContinuumRemainder.sampled_robust_repair_all_centers}\par \nolinkurl{ContinuumRemainder.actual_sampled_routing_blocker_of_stable_miss}\par\smallskip Modules: \nolinkurl{ContinuumGeometric.PeriodicRepair}; \nolinkurl{ContinuumRemainder.RobustRepair}; \nolinkurl{ContinuumRemainder.RobustAssembly} \\\addlinespace
\ref{lem:blocker}. Unconditional uniform robust blocker & \nolinkurl{ContinuumRemainder.robustCompactBlockerSpec_proved}\par\smallskip Modules: \nolinkurl{ContinuumRemainder.FinalProof} \\\addlinespace
\ref{sec:finish}. Countable exhaustion, coefficient normalization and reflected blockers & \nolinkurl{ContinuumRemainder.robustCompactBlockerSpec_open_exhaustion}\par\smallskip Modules: \nolinkurl{ContinuumRemainder.Exhaustion} \\\addlinespace
\ref{thm:main}. Infinitely many distinct values, strict unit measure and nowhere density & \nolinkurl{ContinuumRemainder.powerRemainder_two_sided}\par \nolinkurl{ContinuumRemainder.infinite_tailValuesOutside_of_hits}\par \nolinkurl{ContinuumRemainder.volume_onePeriodic_inter_Icc}\par \nolinkurl{ContinuumRemainder.continuum_power_target}\par\smallskip Modules: \nolinkurl{ContinuumRemainder.DistinctMisses}; \nolinkurl{ContinuumRemainder.TopologyConclusion}; \nolinkurl{ContinuumRemainder.FinalProof} \\\addlinespace
\ref{cor:compact}. Compact restriction with unchanged avoidance quantifiers & \nolinkurl{ContinuumRemainder.compact_power_avoidance}\par\smallskip Modules: \nolinkurl{ContinuumRemainder.FinalProof} \\\addlinespace
\ref{cor:geometric}. Earlier independently proved all-ratio endpoint; current normalization is a traditional corollary & \nolinkurl{ContinuumGeometric.MainTarget}\par \nolinkurl{ContinuumGeometric.geometric_main_target}\par\smallskip Modules: \nolinkurl{ContinuumGeometric.Target}; \nolinkurl{ContinuumGeometric.MainProof} \\\addlinespace
\ref{app:formal}. Exact partial-tail/total-extension correspondence & \nolinkurl{ContinuumRemainder.powerRemainderOn_totalTailExtension}\par \nolinkurl{ContinuumRemainder.tailValuesOutside_totalTailExtension}\par\smallskip Modules: \nolinkurl{ContinuumRemainder.DistinctMisses} \\\addlinespace
\end{longtable}
\normalsize


\begin{thebibliography}{9}
\bibitem{BGKMW}
A. Burgin, S. Goldberg, T. Keleti, C. MacMahon and X. Wang,
\emph{Large sets avoiding infinite arithmetic / geometric progressions},
arXiv:2210.09284v1 (2022), especially Question~1 and Theorem~1.1.
\url{https://arxiv.org/abs/2210.09284v1}.

\bibitem{KP}
M. N. Kolountzakis and E. Papageorgiou,
\emph{Large sets containing no copies of a given infinite sequence},
Analysis \& PDE \textbf{18} (2025), 93--108;
arXiv:2208.02637v2 (2023), especially Theorem~4.1.
\url{https://doi.org/10.2140/apde.2025.18.93}.

\bibitem{FLX}
D.-j. Feng, C.-K. Lai and Y. Xiong,
\emph{Erd\H{o}s similarity problem via bi-Lipschitz embedding},
arXiv:2312.01319v1 (2023), Theorem~1.1.
\url{https://arxiv.org/abs/2312.01319v1}.

\bibitem{Routing}
OpenAI, \emph{The geometric case of the Erd\H{o}s similarity conjecture},
OpenAI/math research family 084, dated 5 October 2026,
research-source snapshot at commit
\texttt{adc7f1241b42e322a6451854ab7e4b4c146bf78a}.
\url{https://github.com/OpenAI/math/tree/adc7f1241b42e322a6451854ab7e4b4c146bf78a}.
Path: \nolinkurl{preprints/The-geometric-case-of-the-Erdos-similarity-conjecture-October-5-2026}.
The corporate attribution follows the source's author field; this is a
research-source citation, not a journal-publication assertion.

\bibitem{Profiles}
mxym, \emph{Modulus-controlled nonlinear avoidance and a differentiability
boundary for null patterns} (version 1), and
\emph{Log-bi-Lipschitz profile avoidance for null configurations} (version 2),
mxym/math entry 006, research-source snapshot at
commit \texttt{87868bcb65bf0460a90d6bc2efc3324e46481230}.
\url{https://github.com/mxym/math/tree/87868bcb65bf0460a90d6bc2efc3324e46481230}.
Path: \nolinkurl{preprints/006-modulus-nonlinear-similarity}.
The repository-account attribution and AI-assisted manuscript status follow
the source; no institutional affiliation or journal publication is asserted.
\end{thebibliography}
\end{document}
